\subsection{Praktische Messung mit dem Oszilloskop}
\label{subsec:rc-schaltung-oszilloskop}

Die theoretisch hergeleitete Aufladekurve einer RC-Schaltung lässt sich mit
einem Funktionsgenerator und einem Oszilloskop sehr gut experimentell
überprüfen. Dazu wird an die RC-Reihenschaltung eine Rechteckspannung angelegt.
Jede steigende Flanke des Rechtecksignals entspricht näherungsweise einem
Spannungssprung von $0$ auf $U_0$.

Der Kondensator lädt sich nach diesem Spannungssprung nicht schlagartig auf,
sondern gemäß der exponentiellen Funktion

\[
U_C(t)=U_0\left(1-e^{-\frac{t-t_0}{\tau}}\right)
\]

mit der Zeitkonstante

\[
\tau=RC.
\]

Bei der fallenden Flanke des Rechtecksignals entlädt sich der Kondensator
wieder exponentiell.

\subsubsection{Geeignete Bauteilwerte}
\label{subsubsec:geeignete-bauteilwerte-rc}

Ein gut messbarer Startwert ist beispielsweise

\[
R = 10\,\text{k}\Omega,
\qquad
C = 100\,\text{nF}.
\]

Damit ergibt sich

\[
\tau = RC
=
10\,000\,\Omega \cdot 100 \cdot 10^{-9}\,\text{F}
=
1\,\text{ms}.
\]

Nach einer Zeitkonstante hat der Kondensator etwa $63{,}2\,\%$ seiner
Endspannung erreicht. Nach etwa fünf Zeitkonstanten gilt der Ladevorgang
praktisch als abgeschlossen.

\[
5\tau = 5\,\text{ms}
\]

Damit sollte die halbe Periodendauer des Rechtecksignals ungefähr

\[
\frac{T}{2} \approx 5\tau
\]

betragen. Also gilt näherungsweise:

\[
T \approx 10\tau.
\]

Für $\tau = 1\,\text{ms}$ folgt:

\[
T \approx 10\,\text{ms}
\]

und damit

\[
f=\frac{1}{T}\approx 100\,\text{Hz}.
\]

Eine sinnvolle Einstellung des Funktionsgenerators ist daher:

\[
f = 100\,\text{Hz}.
\]

\subsubsection{Messaufbau}
\label{subsubsec:messaufbau-rc-oszilloskop}

Der Funktionsgenerator liefert eine Rechteckspannung. Diese wird an die
RC-Reihenschaltung angelegt. Das Oszilloskop misst gleichzeitig die
Eingangsspannung und die Kondensatorspannung.

\[
\text{Generator} \rightarrow R \rightarrow C \rightarrow \text{Masse}
\]

Kanal 1 des Oszilloskops misst die Eingangsspannung $U_0(t)$.
Kanal 2 misst die Kondensatorspannung $U_C(t)$.

\[
\begin{array}{c|c|c}
	\text{Kanal} & \text{Messpunkt} & \text{Signal} \\
	\hline
	\text{CH1} & \text{Generatorausgang} & U_0(t) \\
	\text{CH2} & \text{Kondensator} & U_C(t)
\end{array}
\]

Auf dem Oszilloskop sollte auf Kanal 1 ein Rechtecksignal sichtbar sein.
Auf Kanal 2 erscheint dagegen keine rechteckige Spannung, sondern eine
exponentielle Lade- und Entladekurve.

\subsubsection{Erwartetes Messergebnis}
\label{subsubsec:erwartetes-messergebnis-rc}

Bei einer Rechteckspannung von beispielsweise $0$ bis $5\,\text{V}$ steigt die
Kondensatorspannung nach jeder positiven Flanke exponentiell gegen
$5\,\text{V}$ an:

\[
U_C(t)=5\,\text{V}\left(1-e^{-\frac{t-t_0}{\tau}}\right).
\]

Nach einer Zeitkonstante gilt:

\[
U_C(t_0+\tau)
\approx
0{,}632 \cdot 5\,\text{V}
\approx
3{,}16\,\text{V}.
\]

Nach etwa fünf Zeitkonstanten ist der Kondensator nahezu vollständig geladen:

\[
U_C(t_0+5\tau)\approx 0{,}993\,U_0.
\]

Bei der fallenden Flanke entlädt sich der Kondensator wieder exponentiell:

\[
U_C(t)=U_{C0}e^{-\frac{t-t_0}{\tau}}.
\]

\begin{DidacticBox}[Theorie und Messung]
	\label{box:theorie-und-messung-rc-schaltung}
	Die RC-Schaltung zeigt unmittelbar, wie eine Differentialgleichung in der
	Realität sichtbar wird. Die mathematische Lösung
	
	\[
	U_C(t)=U_0\left(1-e^{-\frac{t-t_0}{\tau}}\right)
	\]
	
	ist nicht nur eine Formel, sondern erscheint direkt als Kurve auf dem
	Oszilloskop. Die Zeitkonstante
	
	\[
	\tau=RC
	\]
	
	kann aus den Bauteilwerten berechnet und anschließend im Messbild überprüft
	werden.
\end{DidacticBox}